% =====================================================================
% 1 INTRODUCTION (house pattern: problem | known | gap | purpose + hypotheses)
% Revised 2026-11-06: H1 = like-for-like contrast (1.x general vs 2.0 DiM); SOM and J-space labelled as
% project-fork plug-ins; 2512.13655 presented as prior work (this is a new paper); positioning sentence added.
% Revised 2026-11-06 (draft v3): hypotheses unchanged. Added (i) the external tools to the design sentence,
% (ii) the KL<=0.1 re-export described as the check of the selection-cap explanation that the earlier draft had
% named as the decisive test (specified after the 32B H1 result, before the re-exported model was evaluated;
% not a new hypothesis), (iii) companion-cell count 13 -> 19 (13 Qwen2.5 cells + Qwen3.5-4B original and its five
% 2.0/fork cells; cells.csv provenance "paper 33" / ref_dir runs/mc-rtx-1), (iv) external tools descriptive only.
% =====================================================================
\section{Introduction}
\label{sec:intro}

Instruction-tuned language models with open weights are released with safety training that makes them refuse
harmful requests, and that refusal can be removed after release without further training. Abliteration edits the
weight matrices that write to the residual stream so that the model can no longer represent a refusal direction
\cite{arditi2024refusal,labonne2024abliteration}, and a growing set of open-source tools automates the edit
\cite{failspy_abliterator,erisforge,jimplus_llmabliteration,heretic2025,obliteratus2026,jwest33abliterator}.
Researchers who study refusal, red-teamers who need unfiltered baselines and auditors who estimate how easily a
release can be stripped of its safeguards must therefore choose a tool, a tool version, a method within the tool
and the rule that selects which edited model is released. The same choice made by a malicious actor determines how
much protection a release retains, so the relative effectiveness of these tools is a safety-relevant quantity
rather than a matter of convenience.

Refusal in many chat models is mediated largely by one direction in activation space \cite{arditi2024refusal}, and
successors of the original difference-in-means (DiM) edit project the direction away from the harmless mean
\cite{lai2024projected}, preserve weight-row magnitudes \cite{lai2025mpoa}, ablate several directions found by
self-organising maps \cite{piras2025som} or replace the direction with an optimised low-rank map
\cite{weidmann2026ara}. Heretic automates the choice of layers and strengths with a Bayesian search that trades
keyword-counted refusals against first-token Kullback--Leibler (KL) divergence \cite{heretic2025,akiba2019optuna}.
Its 2.0 release reorganised the tool around modifier and scorer plug-ins and added arbitrary-rank ablation (ARA)
as a built-in method; the direction-projection and row-normalisation options that define magnitude-preserving
orthogonal ablation (MPOA) were already available in 1.x builds. An earlier comparison \cite{young2025comparative} evaluated
four tools on sixteen 7--14B models and reported large differences in compatibility and in capability
preservation, for example a GSM8K loss of 18.8 percentage points (pp) for Heretic on Yi-1.5-9B against near-zero
losses for single-pass tools.

That comparison shared the measurement practice of the tools themselves. Refusal was counted with keyword markers on
100 prompts from the same AdvBench-derived dataset used to compute the refusal direction, each tool ran at its own
default settings, and distribution shift was summarised by the first-token KL divergence that the search itself
optimises. Keyword matching agrees poorly with human labels and with model-based judges
\cite{xie2025sorrybench,young2026measuring}, and after ablation it can over-count refusals relative to a judge
\cite{young2026harder}; selecting on the same prompts that are reported invites optimistic estimates
\cite{cawley2010overfitting}; and tools that stop at different divergence budgets cannot be ranked without
knowing the budget \cite{dodge2019show}. Version comparisons add a further confound: two releases of the same tool
can apply an identical edit yet be run with different trial budgets, selection rules, system prompts and refusal
markers, so a difference between their outputs need not reflect the version. It is also unknown whether a refusal
direction computed from malicious-code requests, as some practitioners use for coding models, removes more refusal
on such requests than the usual general direction. A practitioner who upgrades from Heretic~1.x to 2.0, runs
another tool at its documented settings or downloads a community abliteration therefore has no held-out,
judge-based evidence about what the choice buys or costs.

The primary purpose of this study was to compare abliteration tools and tool versions on the same base models with
held-out prompts and a validated judge. Archived Heretic~1.x runs with general- and code-derived refusal directions
were rebuilt for six Qwen2.5 models of 7--32 billion parameters and evaluated, together with Heretic~2.0 methods,
two plug-ins of a project fork of Heretic, two further open-source tools run at their documented settings
(OBLITERATUS and jwest33's abliterator, the latter also on the hybrid-attention Qwen3.5-4B), two community releases
and the unedited originals, with one harness for judged refusal, capability and distribution shift. Nineteen cells
(the Heretic~2.0 and fork-plug-in models on Qwen2.5-7B, Qwen2.5-32B and Qwen3.5-4B, their originals and the two
community releases) were first reported in a companion study \cite{young2026harder}, which fixes the tool and
varies model generation; this study fixes the model family and varies tool, version and direction source. Because
Heretic~1.x with the archived settings and Heretic~2.0 DiM apply the same edit, their contrast is the like-for-like
test of the version. \textbf{H1 (version, like-for-like):} on
the same base model, Heretic~2.0 DiM at the companion study's settings would leave less judged residual refusal than
Heretic~1.x with the general direction at the archived settings, without a larger capability loss. The two arms also
differ in the KL cap of the rule that selects the released trial (0.1 versus 0.01), trial budget, search seed,
optimisation system prompt and refusal-marker list, so H1 tests the version together with these settings, not the
edit; MPOA, ARA and the fork plug-ins change the edit as well and are reported as additional configurations. The
search records pointed to the tighter selection cap of the archived runs as a candidate explanation of the
Qwen2.5-32B result, so the trial that the 2.0 selection rule picks from the archived Qwen2.5-32B 1.x general search
was re-exported and evaluated with the same harness. This check was specified after the Qwen2.5-32B result for H1 was known and
before the re-exported model was evaluated; it tests an explanation of the H1 contrast and is not a further
hypothesis. The secondary purposes were to test whether the source of the refusal direction matters and whether
keyword counts change the comparison. \textbf{H2 (prompt domain):} the 1.x code variant would leave less residual refusal than the
general variant on malicious-code prompts and more on other harmful prompts. \textbf{H3 (measurement):}
keyword-matched refusal would differ from judged refusal on the same responses by a tool-dependent amount, so that
keyword counts misstate differences between tools. A further, descriptive expectation was that residual refusal
would decrease as distribution shift increased. The external tools were not the subject of a hypothesis and are
reported descriptively. The hypotheses were formulated after the Heretic~2.0 models had been analysed for the
companion study and were not preregistered; H1 was narrowed from all 2.0 methods to the like-for-like contrast after
an internal review and before the 7B 1.x cells were judged (Section~\ref{sec:deviations}).
